% ---------------------------------------------------------------------------
% content/results.tex -- a hardness section, an algorithm section that
% restates and proves the theorem from the introduction, and a table.
% Sections carry the same anchor/label pair as in the introduction.
% ---------------------------------------------------------------------------

\linkedsection{hardness}{Hardness}\label{sec:hardness}

\begin{theorem}
\label{thm:hardness}\hypertarget{target:thm:hardness}{}%
\textsc{Template Problem} is \NP-hard, even when all processing times are equal.
\end{theorem}

\begin{proof}
By a polynomial-time reduction from \textsc{3-Partition}; see \cite{GJ79} for the source problem.
The construction is omitted in this template.
\end{proof}

\Cref{thm:hardness} rules out a polynomial-time algorithm unless $\PP=\NP$, which motivates the parameterized algorithm of the next section.

\linkedsection{algorithms}{Algorithms}\label{sec:algorithms}

We now restate the main theorem from the introduction. The statement is generated from the same source, so the two cannot drift apart, and it keeps its original number.

\MainTheorem*

\begin{proof}
\Cref{alg:dp} processes the jobs in order of due date and keeps, for every subset of machines that is busy, the best weight seen so far.
There are $2^k$ subsets and $n$ jobs, so the table has $2^k\cdot n$ entries, each computed in constant time from the previous column.
Correctness follows from \Cref{lem:interval}.
\end{proof}

\begin{algorithm}[htb]
\caption{A dynamic program for \textsc{Template Problem}}
\label{alg:dp}
\KwIn{jobs $1,\ldots,n$ sorted by due date; $k$ machines}
\KwOut{the maximum total weight of a feasible schedule}
$T[\emptyset]\gets 0$; $T[S]\gets-\infty$ for every non-empty $S\subseteq\{1,\ldots,k\}$\;
\For{$j\gets 1$ \KwTo $n$}{
    \ForEach{busy set $S$ and machine $i\notin S$ that is free at time $d_j-p_j$}{
        $T[S\cup\{i\}]\gets\max\{T[S\cup\{i\}],\,T[S]+w_j\}$\;
    }
}
\Return $\max_S T[S]$\;
\end{algorithm}

\Cref{tab:experiments} reports running times of \Cref{alg:dp} on random instances; it exists mainly to show a \texttt{booktabs} table with a caption above it, which is the usual convention for tables.

\begin{table}[htb]
    \centering
    \caption{Running time of \Cref{alg:dp} (seconds, mean of ten runs).}
    \label{tab:experiments}
    \begin{tabular}{@{}rrrr@{}}
        \toprule
        $n$ & $k=4$ & $k=8$ & $k=12$ \\
        \midrule
        1\,000  & 0.01 & 0.12 & 1.9 \\
        10\,000 & 0.09 & 1.20 & 19.4 \\
        100\,000 & 0.91 & 12.1 & 197.0 \\
        \bottomrule
    \end{tabular}
\end{table}
